% ===============================================
% ECEV 33400: Stochastic Processes      Fall 2017
% hw_draft.tex
% Template for first draft of homework submission
% ===============================================


% -----------------------------------------------
% The preamble that follows can be ignored. Go on
% down to the section that says "START HERE" 
% -----------------------------------------------

\documentclass{article}

\usepackage[margin=1in]{geometry} 
\usepackage{amsmath,amsthm,amssymb}
\usepackage[sharp]{easylist}

\newcommand{\R}{\mathbb{R}}
\newcommand{\Z}{\mathbb{Z}}
\newcommand{\N}{\mathbb{N}}
\newcommand{\Q}{\mathbb{Q}}
\newcommand{\C}{\mathbb{C}}
\newcommand{\F}{{\mathbb F}}
\newcommand{\Pb}{{\mathbb P}}
\newcommand{\E}{{\mathbb E}}
\newcommand{\Frac}[2]{\left(\frac{#1}{#2}\right)}
\newcommand{\Eval}[2]{\big|_{#1}^{#2}}
\newcommand{\intinf}{\int_{-\infty}^{\infty}}

\newenvironment{theorem}[2][Theorem]{\begin{trivlist}
\item[\hskip \labelsep {\bfseries #1}\hskip \labelsep {\bfseries #2.}]}{\end{trivlist}}
\newenvironment{lemma}[2][Lemma]{\begin{trivlist}
\item[\hskip \labelsep {\bfseries #1}\hskip \labelsep {\bfseries #2.}]}{\end{trivlist}}
\newenvironment{exercise}[2][Exercise]{\begin{trivlist}
\item[\hskip \labelsep {\bfseries #1}\hskip \labelsep {\bfseries #2.}]}{\end{trivlist}}
\newenvironment{problem}[2][Problem]{\begin{trivlist}
\item[\hskip \labelsep {\bfseries #1}\hskip \labelsep {\bfseries #2.}]}{\end{trivlist}}
\newenvironment{question}[2][Question]{\begin{trivlist}
\item[\hskip \labelsep {\bfseries #1}\hskip \labelsep {\bfseries #2.}]}{\end{trivlist}}
\newenvironment{corollary}[2][Corollary]{\begin{trivlist}
\item[\hskip \labelsep {\bfseries #1}\hskip \labelsep {\bfseries #2.}]}{\end{trivlist}}

\newenvironment{solution}{\begin{proof}[Solution]}{\end{proof}}


\begin{document}

% ------------------------------------------ %
%                 START HERE                 %
% ------------------------------------------ %

\title{Homework 3} % Replace X with the appropriate number
\author{ \iffalse STUDENT NAME \fi \\ECEV 33400} % Replace  "\iffalse X \fi" with your name.
\date{Due Date: Thursday 10/26/2017} % Don't change this! :( 

\maketitle

% -----------------------------------------------------
% The following two environments (theorem, proof) are
% where you will enter the statement and proof of your
% first problem for this assignment.
%
% In the theorem environment, you can replace the word
% "theorem" in the \begin and \end commands with
% "exercise", "problem", "lemma", etc., depending on
% what you are submitting. Replace the "x.yz" with the
% appropriate number for your problem.
%
% If your problem does not involve a formal proof, you
% can change the word "proof" in the \begin and \end
% commands with "solution".
% -----------------------------------------------------

\subsection*{Problems} 
% You can copy these problems down to the environments below to write your solutions & delete the top subsections.

\begin{easylist}[enumerate]
# Consider two real-valued random variables, $X_1$ and $X_2$, related to each other through a linear transformation, $X_2 = a\,X_1\,+\,b$. Show that their corresponding probability density functions are related by a multiplicative constant.  

# Consider the stochastic model underlying the theory of island biogeography in the limit of  $S \rightarrow \infty$. Assume a system of $K$ islands along an environmental gradient that makes extinction change from island to island. Therefore, species arrive to the island at the same rate, $\lambda$, the colonization rate, but, they get extinct at a different rates, $\mu_i$, $i=1, \dots, K$. The configuration of the system at any give time is fully defined by the vector $(n_1, \dots, n_K)$, the $K$-dimensional state variable of the system.
## List the total number of events changing system configuration, $(n_1, \dots, n_K)$, at any given time. 
## Calculate the total rate of any configuration change happening in the system at any given time. 
## Using Gillespie algorithm, build stochastic replicates and store the state of the system at increasing times, $t_i= i\,\Delta t$. In different plots, represent the total number of species in the archipelago, and the number of species in an island of intermediate extinction. Start from an empty archipelago and use the following parameters:
### $K=100$, $\lambda = 10.0$, $\mu_i = \mu_0 p^{i-1}$ with $i=1, \dots, K$, $p = 0.99$ and $\mu_0 = 1.0$.
### $K=10000$, $\lambda = 10.0$, $\mu_i = \mu_0 p^{i-1}$ with $i=1, \dots, K$, $p = 0.9999$ and $\mu_0 = 1.0$. Discuss optimal ways of updating the total rate and choosing the next event to happen. 
## What is the distribution for the total number of species in the archipelago, $N = \sum_{i=1}^K n_i$, at stationarity? Through simulation, assess that the distribution of this number over your stochastic replicates tends to this theoretical distribution as time increases.  
\end{easylist}

% -----------------------------------------------------
% First problem
% -----------------------------------------------------

% \begin{problem}{1}
% Show that the probability-generating function for the binomial distribution $p_j = \binom{n}{j} p^j q^{n-j}$, where $ q = 1 - p$, is $P(x) = (q + px)^n$. Hence prove that the mean is $m = np$, and the variance is $\sigma^2 = npq$.
% \end{problem}

% \begin{solution}
% \end{solution}

% -----------------------------------------------------
% Second problem
% -----------------------------------------------------

% \vspace{0.25in} % This just adds some space between problems 1 and 2.

% \begin{problem}{2}
% What is the probability-generating function for the Poisson distribution $$p_j = \frac{e^{-\lambda} \lambda^j}{j!}, j = 0, 1, 2, \cdots ?$$ Hence show that the mean and variance of the distribution are both equal to $\lambda$. 
% \end{problem}

% \begin{solution}
% \end{solution}

% -----------------------------------------------------
% Third problem
% -----------------------------------------------------


% -----------------------------------------------------
% Fourth problem
% -----------------------------------------------------


% -----------------------------------------------
% Ignore everything that appears below this.
% -----------------------------------------------

\end{document}
